\section{Conclusions}
\label{sec:conclusion}

This paper studies residual connectivity for graph classification as a structured branch-by-layer design problem. Across six datasets and four message-passing operators, residual topology matters, but its value depends on the dataset, operator, and branch budget. MatrixRes is the most frequent winner and has the best mean rank in the full benchmark, but the per-dataset GCNConv winners remain mixed: Plain is strongest on PROTEINS, HorizontalRes on DD, MatrixResGated on ENZYMES, and MatrixRes on MUTAG, AIDS, and Mutagenicity. The ENZYMES comparison has low absolute accuracy and should be interpreted cautiously. The branch-count ablations further show that adding branches is useful only within a limited operating region.

The mechanism analysis provides a plausible explanation for why branch expansion is not automatically beneficial. Additional branches can increase representational diversity, but they also increase residual traffic and may preserve redundant branch similarity or weaken optimization. Matrix-style reuse is best viewed as a controlled residual family whose effectiveness depends on the branch budget and residual filtering strategy.

The practical implication is that residual topology, branch count, operator choice, and residual control should be evaluated jointly. Before treating a sensitivity result as a performance claim, candidate settings should be rerun across folds. In the current experiments, this validation supports a smaller and more parameter-efficient MatrixResGated model on DD, while the PROTEINS sparsity candidate remains less stable.
